\section{Foundations of Digital Forensics}

Digital forensics involves the collection, analysis, and preservation of digital evidence for legal purposes. However, the process is complicated by challenges in legislation, jurisdiction, and technical limitations such as encryption and data retention laws. The collaboration between law enforcement and service providers, such as Google, is a central part of the investigation process.

\subsection{Types of Digital Evidence}
Digital evidence can be classified into three categories:
\begin{itemize}
    \item \textbf{Created by Man:} Data generated through human interaction, such as emails or documents
    \item \textbf{Created Independently by Computers:} Data generated by machines through processing, without human intervention, like telephone records or server logs
    \item \textbf{Created by Both Man and Computer:} Data where human input and computer processing are both involved, such as in a spreadsheet
\end{itemize}

\subsubsection*{Digital Evidence Requirements}
For digital evidence to be admissible in court, it must be:
\begin{itemize}
    \item \textbf{Admissible:} Compliant with the law and best practices.
    \item \textbf{Authentic:} Free from tampering.
    \item \textbf{Reliable and Believable:} Understandable to a judge.
    \item \textbf{Proportional:} Respectful of the rights of all parties involved.
\end{itemize}
These principles ensure the integrity and value of the evidence in court.

\subsection{Legal and Jurisdictional Issues in Digital Forensics}
The complexity of digital forensics is compounded by jurisdictional challenges. \\
Laws governing data retention and privacy vary across countries, leading to conflicts between international legal systems.
\smallskip

\noindent
For example, in the EU, data retention directives mandate retention periods of up to 24 months, but this law does not apply to US-based companies like Google, which comply voluntarily with European law.

\subsubsection*{Data Retention and Voluntary Compliance}
Google, based in the US, has entities in multiple countries, including Ireland. Under the GDPR, Google Ireland is responsible for handling data for European users. However, there is no obligation under European law for Google to retain data. \\
This is a voluntary practice that provides challenges for law enforcement, as they must rely on the company's willingness to comply with requests.

\subsection{Encryption and Remote Forensics}
Encryption is one of the greatest obstacles in digital forensics, as it prevents access to data on suspect devices. However, some propose the use of remote forensics, where law enforcement uses spyware or keyloggers to access encrypted data on a suspect's device. This approach is controversial, as it involves bypassing encryption without the user's consent.

\subsubsection*{The Ethical Dilemma of Remote Forensics}
Remote forensics introduces ethical questions about the extent to which law enforcement should be allowed to access personal data. The use of spyware and trojans to hack into a device without a warrant raises concerns about privacy and the potential for abuse of power.

\subsubsection*{Key Mandatory Disclosure Laws and Their Limitations}
Key mandatory disclosure laws require suspects to provide encryption keys, but their effectiveness is limited by the possibility of anti-forensics techniques, where criminals hide or encrypt data within layers of encryption. Additionally, such laws may violate human rights, particularly the right to remain silent as outlined in Article 6 of the European Convention on Human Rights.

\subsubsection*{Remote Forensics}
Remote forensics, it is a hacking into a suspect's device without their consent. It has been proposed as a solution to bypass encryption. However, this raises significant legal and ethical issues.\\
The German Constitutional Court, for example, ruled against laws allowing law enforcement to hack into computers, citing violations of privacy and the right to informational self-determination.

\subsubsection{Cloud Computing: Jurisdiction and Data Location Issues}
Cloud computing presents unique challenges for digital forensics due to the "loss of location" of data. In a cloud environment, data can be stored across multiple physical locations, and countries, making it difficult to determine which jurisdiction applies when law enforcement seeks access to evidence.
\bigskip

\noindent
The power of disposal approach allows law enforcement to access data stored in the cloud if they have control over the device. This approach circumvents traditional jurisdictional principles but raises concerns about the balance between security and privacy rights

\subsection{Future of Digital Forensics}
As digital forensics continues to evolve, new technologies and methods must be developed to address the challenges posed by encryption, cloud computing, and jurisdictional issues. Collaboration between international law enforcement agencies and service providers will be essential to ensure that justice is served while protecting individuals' privacy rights.

\begin{table}[H]
\centering
\small
\begin{tabular}{|p{3.5cm}|p{5cm}|p{3.2cm}|}
\hline
\textbf{Jurisdiction Principle} & \textbf{Description} & \textbf{Example} \\
\hline
Territorial Principle & Jurisdiction based on the physical location of data & Data stored in a US server is governed by US law. \\
Nationality Principle & Jurisdiction based on the nationality of the perpetrator & A Russian citizen commits a crime in Italy. \\
Flag Principle & Jurisdiction based on the location of a vessel or aircraft & A crime committed on a US-flagged ship is governed by US law. \\
Power of Disposal Principle & Jurisdiction based on control over data & Law enforcement can access data stored anywhere if they control the device. \\
\hline
\end{tabular}
\caption{Different Jurisdiction Principles in Digital Forensics.}
\end{table}

\subsection{The Budapest Convention on Cybercrime}

In 2001 the \textbf{Budapest Convention on Cybercrime} was adopted as the first international treaty addressing crimes committed through computer systems and networks.\\
The convention was developed within the framework of the \textbf{Council of Europe} (an international organization based in Strasbourg that operates independently from the institutions of the European Union).
\medskip

\noindent
The objective of the convention is to establish a common legal framework for combating cybercrime and to facilitate cooperation between states during digital investigations.\\
Key characteristics of the convention include:

\begin{itemize}
\item \textbf{Harmonization of legislation}: defining common categories of cybercrime that participating countries should criminalize within their national legal systems
\item \textbf{Procedural powers}: providing investigative tools and legal procedures specifically designed for digital evidence and computer systems
\item \textbf{International cooperation}: enabling cross-border collaboration between law enforcement authorities during cybercrime investigations
\end{itemize}

\medskip

\noindent
The convention was drafted in a technological context very different from today. At the beginning of the 2000s:

\begin{itemize}
\item personal computers were only beginning to spread within households;
\item most digital infrastructures were based on \textit{on-premise} systems;
\item cloud computing services were not yet widely adopted.
\end{itemize}

\medskip

\noindent
Despite this context, the convention anticipated many of the challenges that later became central in cybercrime investigations, particularly those related to jurisdiction and international cooperation.

\subsubsection *{Signature and Ratification of International Conventions}

International conventions generally involve two distinct phases:

\begin{itemize}

    \item \textbf{Signature}: the signature of a convention indicates that a state agrees with the general principles and objectives of the treaty. However, the signature alone does not make the convention legally binding within the national legal system.

    \item \textbf{Ratification}: ratification represents the formal adoption of the convention within the domestic legal framework of a country. This process usually requires the approval of national legislative bodies and the implementation of laws that comply with the provisions of the convention.

\end{itemize}
\noindent
Without ratification, a convention remains primarily a political commitment rather than a legally enforceable instrument. Therefore, the effectiveness of international agreements depends largely on whether states proceed with the ratification process and integrate the convention into their national legal systems.
\smallskip

\noindent
The time required for ratification can vary significantly between countries. For example, Italy signed the Budapest Convention in 2001 but completed the ratification process only in 2008. This delay reflects the complexity of adapting national legislation to comply with international agreements.



\subsubsection{International Criticism and Opposition}

Despite its wide adoption, the convention has also faced criticism and political opposition from several countries.

\medskip
\noindent
Main concerns raised by non-participating states include:

\begin{itemize}
    \item \textbf{Sovereignty concerns}: some countries argued that certain provisions could limit national control over digital infrastructures and investigations

    \item \textbf{Data-sharing concerns}: governments expressed reservations regarding the exchange of digital evidence with foreign law enforcement agencies

    \item \textbf{Lack of participation in drafting}: some states, such as India, initially criticized the convention because they were not involved in the drafting process
\end{itemize}

\noindent
In addition, countries such as Russia ultimately rejected the convention, citing concerns related to sovereignty and international investigative cooperation.



\subsubsection{Crimes Addressed by the Convention}

The Cybercrime Convention focuses on crimes that involve the use of digital technologies either as a target or as a tool.
\medskip

\noindent
Two main categories of offences are considered:

\begin{itemize}
    \item \textbf{Crimes against computer systems}: offences where the digital device or infrastructure itself is the target of the attack.

    \item \textbf{Technology-enabled crimes}: traditional crimes that are facilitated or enhanced through the use of digital technologies, such as fraud or financial crimes.
\end{itemize}

\noindent
The convention therefore provides guidelines for states developing national legislation on cybercrime and establishes a framework for international cooperation in investigations involving digital evidence.


\subsubsection*{International Cooperation Principles}

A central component of the convention is the establishment of mechanisms that facilitate cooperation between countries during cybercrime investigations.\\
Key cooperation mechanisms include:

\begin{itemize}
    \item \textbf{Mutual assistance}: enabling states to request investigative support or digital evidence from other countries.

    \item \textbf{Expedited preservation of data}: allowing authorities to rapidly preserve stored computer data that may otherwise be deleted or modified.

    \item \textbf{Real-time collection of traffic data}: enabling authorities to monitor communication metadata during investigations.

    \item \textbf{Interception of content data}: allowing the interception of communication content under specific legal conditions.
\end{itemize}

\noindent
To improve the efficiency of international cooperation, the convention requires each participating state to designate a contact point available at all times.\\
Key characteristics include:

\begin{itemize}
    \item availability of a national contact point \textbf{24 hours a day, 7 days a week};

    \item rapid exchange of information related to cybercrime investigations;

    \item support for urgent requests involving the preservation or collection of electronic evidence.
\end{itemize}

\noindent
This mechanism was inspired by the G8 network of cybercrime contact points and aims to reduce delays caused by traditional bureaucratic procedures in international legal cooperation.


\subsection{Articles}

\subsubsection*{Article 14: Scope of Procedural Provisions}

Article 14 defines the scope of the procedural powers established by the Convention.\\
Each party must adopt legislative measures enabling authorities to investigate:

\begin{itemize}
\item offences established in accordance with the Convention;
\item other offences committed through a computer system;
\item any criminal offence involving electronic evidence.
\end{itemize}



\subsubsection*{Article 15: Conditions and Safeguards}

Article 15 establishes the fundamental safeguards governing the application of investigative powers.

\begin{itemize}
\item investigative measures must respect human rights and fundamental freedoms;
\item authorities must ensure proportionality of the investigative action;
\item judicial or independent supervision must be guaranteed;
\item the rights of third parties must be taken into account.
\end{itemize}



\subsubsection*{Article 16: Expedited Preservation of Stored Computer Data}

Article 16 introduces the possibility for authorities to require the rapid preservation of stored computer data.

\begin{itemize}
\item authorities may order a person or service provider to preserve specific data;
\item preservation is used when data is at risk of deletion or modification;
\item the preservation order generally applies for up to 90 days and may be extended.
\end{itemize}



\subsubsection*{Article 17: Expedited Preservation and Disclosure of Traffic Data}

Article 17 extends the preservation mechanism to traffic data.

\begin{itemize}
\item authorities may require the rapid preservation of traffic data related to a specific communication;
\item preserved traffic data can be used to identify the service providers involved;
\item the measure allows investigators to reconstruct the communication pathway.
\end{itemize}



\subsubsection*{Article 18: Production Order}

Article 18 allows authorities to order the disclosure of specific computer data.

\begin{itemize}
\item a person located within the territory may be ordered to provide computer data in their possession;
\item service providers may be required to disclose subscriber information;
\item the measure applies only to entities located within the requesting state's jurisdiction.
\end{itemize}



\subsubsection*{Article 19: Search and Seizure of Stored Computer Data}

Article 19 provides authorities with powers to access and secure computer data during investigations.

\begin{itemize}
\item authorities may search computer systems located within their territory;
\item stored computer data may be accessed, copied, or seized;
\item storage media containing digital evidence may also be secured.
\end{itemize}



\subsubsection*{Article 20: Real-Time Collection of Traffic Data}

Article 20 allows authorities to collect traffic data in real time during an investigation.

\begin{itemize}
\item traffic data may be collected directly by law enforcement authorities;
\item service providers may be required to assist in the collection process;
\item the measure enables monitoring of communication metadata.
\end{itemize}



\subsubsection*{Article 21: Interception of Content Data}

Article 21 allows the interception of communication content in real time for serious offences.

\begin{itemize}
\item authorities may intercept the content of communications;
\item interception may be carried out directly or through service providers;
\item the measure must comply with national legal safeguards.
\end{itemize}



\subsubsection*{Article 29: Expedited Preservation of Stored Computer Data (Mutual Assistance)}

Article 29 allows a party to request another party to preserve stored computer data located in its territory.

\begin{itemize}
\item preservation ensures that data remains available for future investigative requests;
\item the request is typically followed by a formal mutual assistance procedure;
\item the objective is to prevent deletion of evidence during international investigations.
\end{itemize}



\subsubsection*{Article 30: Expedited Disclosure of Preserved Traffic Data}

Article 30 allows the disclosure of preserved traffic data between parties.

\begin{itemize}
\item the requested party must disclose sufficient traffic data to identify service providers involved;
\item the information helps determine the transmission path of the communication;
\item this mechanism supports cross-border investigations.
\end{itemize}



\subsubsection*{Article 32: Trans-Border Access to Stored Computer Data (Most important article for Vaciago)}

Article 32 regulates the possibility of accessing data located in another state without formal mutual assistance procedures.

\begin{itemize}
\item \textbf{Article 32(a)} allows access to publicly available data regardless of the physical location of the data.

\item \textbf{Article 32(b)} allows access to data located in another state if the investigator obtains the lawful and voluntary consent of the person authorized to disclose the data.
\end{itemize}

\noindent
Article 32(b) is particularly controversial because the disclosure of data ultimately depends on the voluntary cooperation of the entity controlling the data, typically the service provider.